\documentclass[11pt,a4paper]{article}
\usepackage[top=0.65in,bottom=0.65in,left=0.75in,right=0.75in]{geometry}
\usepackage{fontspec}
\usepackage[english,hindi,provide=*]{babel}
\babelprovide[onchar=ids fonts]{english}
\babelfont[hindi]{rm}[Renderer=HarfBuzz,Script=Devanagari,BoldFont={Noto Sans Devanagari Bold},ItalicFont={Noto Sans Devanagari},ItalicFeatures={FakeSlant=0.15},BoldItalicFont={Noto Sans Devanagari Bold},BoldItalicFeatures={FakeSlant=0.15}]{Noto Sans Devanagari}
\babelfont[english]{rm}{Latin Modern Roman}
\renewcommand{\bfdefault}{b}
\usepackage{enumitem}
\usepackage{hyperref}
\usepackage{titlesec}

\setlist{nosep,leftmargin=*}
\setlist[itemize]{label=\foreignlanguage{english}{\textbullet}}
\pagestyle{empty}
\hypersetup{colorlinks=true,urlcolor=black}
\titleformat{\section}{\large\bfseries}{}{0em}{}[\rule{\textwidth}{0.4pt}]
\titlespacing{\section}{0pt}{7pt}{3pt}

\begin{document}

\begin{center}
  {\LARGE\textbf{आरव शर्मा}}\\[3pt]
  वरिष्ठ सॉफ्टवेयर अभियंता\\[2pt]
  बेंगलुरु, भारत \quad|\quad
  \href{mailto:aarav.sharma@example.com}{aarav.sharma@example.com} \quad|\quad
  +91 98765 43210
\end{center}

\section{व्यावसायिक सारांश}
आठ वर्षों के अनुभव वाले सॉफ्टवेयर अभियंता, जिन्होंने विश्वसनीय वितरित प्रणालियाँ बनाई हैं और अलग-अलग टीमों के साथ बड़े तकनीकी कार्यक्रम सफलतापूर्वक पूरे किए हैं।

\section{कार्य अनुभव}
\noindent\textbf{वरिष्ठ सॉफ्टवेयर अभियंता} \hfill \textit{जुलाई 2021 -- वर्तमान}\\
\textit{उदाहरण टेक्नोलॉजी, बेंगलुरु}
\begin{itemize}
  \item भुगतान मंच की नई सेवा तैयार की, जिससे प्रतिदिन दस लाख अनुरोध सुरक्षित रूप से संसाधित हुए और औसत विलंब तीस प्रतिशत घटा।
  \item छह अभियंताओं की टीम का मार्गदर्शन किया और स्वचालित परीक्षण बढ़ाकर उत्पादन त्रुटियों में चालीस प्रतिशत कमी की।
  \item उत्पाद और सुरक्षा टीमों के साथ मिलकर चरणबद्ध माइग्रेशन पूरा किया तथा सभी ग्राहकों के लिए सेवा लगातार उपलब्ध रखी।
\end{itemize}

\noindent\textbf{सॉफ्टवेयर अभियंता} \hfill \textit{जून 2018 -- जून 2021}\\
\textit{नमूना सिस्टम्स, पुणे}
\begin{itemize}
  \item निगरानी उपकरण विकसित किए जिनसे समस्याओं की पहचान तेज हुई और सहायता दल का साप्ताहिक समय बचा।
  \item डेटा प्रसंस्करण कार्यप्रवाह सुधारा और मासिक आधारभूत ढाँचा लागत पंद्रह प्रतिशत कम की।
\end{itemize}

\section{शिक्षा}
\textbf{कंप्यूटर विज्ञान में प्रौद्योगिकी स्नातक} \hfill \textit{2018}\\
भारतीय प्रौद्योगिकी संस्थान, दिल्ली

\section{तकनीकी कौशल}
पाइथन, टाइपस्क्रिप्ट, पोस्टग्रेएसक्यूएल, रेडिस, डॉकर, कुबेरनेटिस, एडब्ल्यूएस

\section{भाषाएँ}
हिन्दी — मातृभाषा \quad|\quad अंग्रेज़ी — व्यावसायिक दक्षता

\end{document}
